\documentclass[11pt]{article}
\usepackage[margin=1.1in]{geometry}
\usepackage{amsmath,amssymb,amsthm,mathtools}
\usepackage[T1]{fontenc}
\usepackage{microtype}
\usepackage{booktabs}
\usepackage{graphicx}
\usepackage{xcolor}
\usepackage{algorithm}
\usepackage{algpseudocode}
\usepackage{enumitem}
\usepackage[colorlinks=true,linkcolor=black,citecolor=black,urlcolor=black]{hyperref}

\newcommand{\R}{\mathbb R}
\newcommand{\E}{\mathbb E}
\newcommand{\Q}{\mathbb Q}
\newcommand{\PP}{\mathbb P}
\newcommand{\dd}{\mathrm d}
\newcommand{\dlv}{\Sigma}
\newcommand{\abs}[1]{\lvert #1\rvert}
\newcommand{\ip}[2]{\langle #1,#2\rangle}
\newcommand{\qv}[1]{\langle #1\rangle}
\newcommand{\one}{\mathbf 1}
\newcommand{\Var}{\operatorname{Var}}
\newcommand{\Gh}{\widehat G}
\newcommand{\Nh}{\widehat N}
\newcommand{\removeblock}[1]{}
\newcommand{\MW}[1]{{\color{red}#1}}
\DeclareMathOperator{\diag}{diag}
\DeclareMathOperator{\tr}{tr}

\theoremstyle{plain}
\newtheorem{theorem}{Theorem}[section]
\newtheorem{proposition}[theorem]{Proposition}
\newtheorem{lemma}[theorem]{Lemma}
\newtheorem{corollary}[theorem]{Corollary}
\theoremstyle{definition}
\newtheorem{assumption}[theorem]{Assumption}
\theoremstyle{remark}
\newtheorem{remark}[theorem]{Remark}

\title{Universal Approximation Theorem\\{Work in Progress}}
\author{}
\begin{document}
\maketitle

We want to derive a UAT that GROMM can approximate any LVSV model for a finite number of maturities. It turns out that the first part of this is an alternate derivation of the BHW drift restriction. 

Assume for a fixed list of strikes and maturities the model prices are defined as $C = (C^{j, i})$ for strikes $K^0 < \dots < K^{N+1}$ and expirations $T_1 < \dots < T_M$. Since prices are risk-neutral under $\Q$ we know in particular that for a SV model, any call price $C$ is a positive martingale (stochastic exponential), i.e. $d C^{j, i}_t = C^{j, i}_t\psi^{j, i}_tdW_t $ for some process $\psi^{j, i}_t$. 

Recall that the squared discrete local volatility (DLV squared) is given by 
\begin{equation*}
    \Sigma^{j, i} = \dfrac{b\Theta^{j, i}}{\frac{1}{2}(K^{i})^2\Gamma^{j, i}}
\end{equation*}
where backward Theta and Gamma are defined as: 
\begin{align*}
    b\Theta^{j, i} &= \dfrac{C^{j, i} - C^{j-1, i}}{T_{j} - T_{j-1}}
    \\
    \Gamma^{j, i} &= \dfrac{\frac{C^{j, i + 1} - C^{j, i}}{K^{i+1} - K^i} - \frac{C^{j, i} - C^{j, i-1}}{K^{i} - K^{i-1}}}{\frac{K^{i+1} - K^{i-1}}{2}}
\end{align*}
Observe that Gamma and the risk-neutral density are directly related: 
\begin{equation*}
    q^{j, i} = \frac{K^{i+1} - K^{i-1}}{2} \Gamma^{j, i}. 
\end{equation*}
Write $dq^{j, i}_t = E^{j, i}_t dW $ for $E^{j, i}_t = q^{j, i}_t \eta^{j, i}_t$. 


\begin{equation*}
    d \Sigma^{j, i} / \Sigma^{j, i} = \mu^{j, i} dt + \xi^{j, i} dW 
\end{equation*}

\begin{equation}
    \label{eq:dlnDLV}
    d \ln \Sigma^{j, i} = d \Sigma^{j, i} / \Sigma^{j, i} - \tfrac{1}{2}  d \qv{\Sigma^{j, i}} / (\Sigma^{j, i})^2 = \mu^{j, i} dt + \xi^{j, i} dW - \frac{1}{2} |\xi^{j, i}|^2 dt 
\end{equation}


We apply Itô on $\ln \Sigma^{j, i} = \ln b \Theta^{j, i} - \ln q^{j, i} + \text{const}$. Suppress dependence on current time $t$ for readability. 
\begin{equation*}
    d \ln b \Theta^{j, i} =  \zeta^{j, i} dW - \frac{1}{2} |\zeta^{j, i}|^2 dt  \quad \textrm{for} \quad \zeta^{j, i} \coloneqq \dfrac{C^{j, i}\psi^{j, i} - C^{j-1, i} \psi^{j-1, i}}{C^{j, i} - C^{j-1, i}}
\end{equation*}

\begin{equation*}
    d \ln q^{j, i} = \eta^{j, i} dW - \frac{1}{2} |\eta^{j, i}|^2 dt  
\end{equation*}

Combining the expressions we obtain Itô on the log-DLV squared: 
\begin{equation}
    \label{eq:intermediate}
    d \ln \Sigma^{j, i} = \frac{1}{2} \left( |\eta^{j, i}|^2 -  |\zeta^{j, i}|^2 \right)dt + \left( \zeta^{j, i}  - \eta^{j, i} \right)dW  
\end{equation}
where $\xi^{j, i} = \zeta^{j, i} - \eta^{j, i}$. Substituting $\zeta^{j, i} = \xi^{j, i} + \eta^{j, i}$ into \eqref{eq:intermediate} gives 
\begin{align*}
    d \ln \Sigma^{j, i} &= \frac{1}{2} \left( |\eta^{j, i}|^2 -  |\eta^{j, i} + \xi^{j, i}|^2 \right)dt + \xi^{j, i} dW  
    \\ 
    &= \frac{1}{2} \left( - 2\ip{\xi^{j, i}}{\eta^{j, i}} - |\xi^{j,i}|^2  \right)dt + \xi^{j, i} dW 
\end{align*}
Compare this with \eqref{eq:dlnDLV} and match the terms to obtain $\mu^{j, i} = -\ip{\xi^{j,i}}{\eta^{j,i}}$, which demonstrates that for any lattice of options $(C^{j, i})_{j, i}$ from a diffusion model the drift coincides with the BHW-type drift restriction. 

Next, we inspect whether the assumptions of the existence theorem fit the assumptions of the SV model and how we can extend them. We have a couple: 
\begin{itemize}
    \item Spot stays in $[K^0, K^{N+1}]$: violation in general for SV models. $K^0 = 0$ works. $K^{N+1} \rightarrow \infty $ is problematic as it's an absorbing state: once we hit it, spot stays there. In an SV model the spot continues. However, we can keep on increasing $K^{N+1}$ such that the hitting probability goes towards zero. But then the representation holds only in the limit, i.e. SV model is a GROMM model with growing top barrier. 
    \item Lipschitz continuity of $\xi$: in general not true. Volatility process $\sqrt V$ in SV model can be Hölder-$\tfrac{1}{2}$ (Heston). 
    \item $\xi$ is square-integrable: $\xi = \zeta - \eta$, both $\zeta$ and $\eta$ could explode. Difference needs to be  square-integrable. 
    \item The Gaussian bridge front / settlement on the strike: we need to drop this and just use $q^1$ as implied by the SV model. $q^1$ moves freely but needs to lead to a barycentric payoff on two neighbouring nodes (as determined by the SV model). Extreme nodes stay absorbing (as determined by $\Omega_j$). 
    \item Mass on interior strikes $q^1_t$ at $t=T_{k-1}$: this holds in general for SV models. 
\end{itemize}


\end{document}